\subsection{Detalle de tecnologías seleccionadas}

A continuación se resume la tecnología de cada capa y su justificación principal:

\begin{itemize}
  \item \textbf{Backend (12 módulos).} Laravel 13 sobre PHP 8.5, con Composer y el archivo de versiones de dependencias, organizado en contextos M1--M12. PostgreSQL/PostGIS conserva los datos geográficos; M4 usa repositorios espaciales con SQL parametrizado vía PDO/Query Builder y pruebas de consultas. Laravel 13 requiere PHP 8.3 o superior y admite PHP 8.5 (Laravel, 2026); ambos componentes se actualizarán a versiones con soporte durante los 56 meses.
  \item \textbf{Frontend web.} Angular 22 con Tailwind CSS y TypeScript 6.0 compatible con su matriz oficial. La integración con Keycloak se implementa mediante OIDC y se mantiene junto con las dependencias del cliente. Las versiones deben actualizarse durante el contrato conforme a sus ciclos de soporte.
  \item \textbf{Apps de campo (preventa y reparto).} Kotlin Android nativo. Nativo del parque Zebra/Android, escáner GS1 vía Zebra DataWedge, SQLite/Room offline, acceso nativo a GPS, POS y térmica.
  \item \textbf{Borde y continuidad del acceso (Capa 2).} CloudFront, AWS WAF, AWS Shield Advanced y ALB en nube, con control de acceso local. En los CD, Starlink fijo permanece encendido con IPsec establecido y BGP de menor preferencia tras fibra y LTE; en los cross-docking es el camino principal, con LTE de dos proveedores como respaldo y conmutación automática en menos de 30 segundos (ADR-02), sin sustituir la autonomía del servicio local.
  \item \textbf{Puerta de enlace de servicios (Capa 3).} Amazon API Gateway para publicar las APIs y aplicar controles de entrada. La configuración debe cubrir la validación de identidad de Keycloak, las cuotas y los límites de solicitudes, conforme a la modalidad de API seleccionada y al flujo descrito en la capa de puerta de enlace de este apartado.
  \item \textbf{Base de datos transaccional on-premise.} PostgreSQL + PostGIS por sitio con el mismo perfil local de bodega; cada sitio es autoridad de sus movimientos y N-04/N-05 consolida stock y reservas en nube.
  \item \textbf{Base de datos nube (OLTP y recuperación).} Amazon Aurora PostgreSQL, con despliegue Multi-AZ y una ventana propuesta de recuperación a un punto en el tiempo (PITR) de 35 días. Los tiempos de conmutación y restauración se verifican mediante pruebas frente a los objetivos RTO/RPO establecidos.
  \item \textbf{Ingesta IoT / frío.} Amazon DynamoDB e IoT Core reciben las lecturas; Greengrass corre en tres gateways Moxa UC-8200 o equivalentes: dos en Talca, que leen ambas cámaras por Modbus TCP para evitar un punto único de falla, y uno en Concepción; leen sensores por Modbus, bloquean el despacho ante una excursión crítica y sostenida en menos de 5 segundos y conservan 24 horas sin enlace. Los termógrafos de los camiones registran toda la ruta en memoria interna y alertan por BLE al terminal, que reenvía el aviso al recuperar cobertura y descarga el registro completo al volver. Escrituras serverless con TTL nativo (raw 30 días).
  \item \textbf{Serie consolidada (OLAP).} S3 Parquet + AWS Glue + Redshift Serverless. OLAP por diseño: DynamoDB raw $\rightarrow$ Glue $\rightarrow$ S3 Parquet $\rightarrow$ Redshift, sin motor de series adicional.
  \item \textbf{Caché y datos de turno.} Redis acelera las lecturas centrales; no se exige Redis en cada sitio. Las APIs entregan una copia cifrada para SQLite/Room y conservan por separado la cola de escrituras sin confirmar. La identidad usa credenciales de 8 horas en bodega y 14 horas en terreno; su caché local de solo lectura tiene TTL de 24 horas en VM-05, VM-C03 y cada E-01, junto al verificador local de relevos (ADR-06).
  \item \textbf{Mensajería asíncrona (Capa 5).} RabbitMQ en on-premise mediante adaptador AMQP \texttt{php-amqplib}; SQS FIFO con sobre JSON y consumidor PHP para reconciliación, colas SQS separadas para trabajos Laravel, y SNS para difusión y alertas en nube. Un único planificador Laravel programa las tareas periódicas. \textbf{Dos procesos del servicio de integración con el ERP en VM-04 consumen RabbitMQ local y, por salida, la cola FIFO de solicitudes al ERP; ERP integrado exclusivamente vía ACL; Hub EDI GS1 centralizado en nube (EANCOM, GS1 XML, EPCIS), AS2 como transporte.}
  \item \textbf{Analítica / BI.} S3 Data Lake, Redshift Serverless y Glue ETL. Consultas históricas sin degradar el procesamiento transaccional.
  \item \textbf{Objetos / documentos.} Amazon S3 con Intelligent-Tiering y Object Lock.
  \item \textbf{IAM / Identidad (Capa 7).} Keycloak (OIDC, SAML, SSO, MFA) como IdP maestro en nube, con caché local de solo lectura de TTL 24 horas y verificador local con manifiesto firmado y PIN personal en VM-05, VM-C03 y cada E-01.
  \item \textbf{Gestión de secretos (Capa 7).} AWS Secrets Manager y SSM Parameter Store con rotación automática. Cuenta de emergencia fuera de banda, con doble autorización, registro de uso y rotación posterior de credenciales.
  \item \textbf{Observabilidad (Capa 8).} OpenTelemetry para PHP/Laravel y colectores ADOT on-premise con almacenamiento temporal en disco de 24 horas; plataforma única Amazon CloudWatch para logs (12 meses en línea + 24 en archivo), métricas (13 meses), trazas y tableros (ADR-14).
  \item \textbf{Gestión de dispositivos (RT-03.18; Bases Técnicas Transversales, cap. 3, p. 9).} Gestión centralizada mediante MDM (Android Enterprise / Zebra DNA, SaaS), con inventario, configuración, actualización de firmware y aplicaciones, bloqueo y borrado remoto del parque móvil.
  \item \textbf{Contenedores / orquestación.} Imagen PHP 8.5 con servidor HTTP/PHP-FPM para APIs y procesos PHP CLI separados para colas y programación. ECS Fargate aloja los perfiles centrales; Talca y Concepción ejecutan el perfil local de bodega sobre máquinas virtuales Proxmox, y Docker Compose ejecuta el mismo perfil local de bodega en los equipos E-01 de los cross-docking. El código y las migraciones de esquema proceden del mismo artefacto versionado.
  \item \textbf{Infraestructura como Código.} Terraform administra recursos mediante proveedores y estados separados por ambiente; Ansible configura hosts. CDK no administra los mismos recursos: cada activo tiene un único propietario de infraestructura como código.
  \item \textbf{CI/CD.} GitLab CI como orquestador + AWS CodeBuild para construcción hermética con procedencia SLSA 3; la auditoría de dependencias, PHPUnit, PHPStan/Larastan y Laravel Pint verifican dependencias, comportamiento, análisis estático y formato. \texttt{swagger-php} genera OpenAPI 3.1 desde atributos PHP; los esquemas AsyncAPI 2.6 se validan y publican en la misma puerta de calidad.
\end{itemize}

El Anexo 4-P reúne el inventario completo de servicios AWS y las equivalencias entre nombres funcionales y técnicos. La descripción por capa explica la responsabilidad de cada tecnología; el inventario permite comprobar su correspondencia con el despliegue.

Laravel conserva una superficie de APIs, políticas y colas coherente para los doce módulos;

\begin{itemize}
  \item Django también permitiría el monolito. Se selecciona un único entorno Laravel/PHP para evitar operar dos cadenas de dependencias y dos familias de trabajadores.
  \item La equivalencia se acepta por contratos y pruebas, no por parecido de bibliotecas.
  \item PostgreSQL/PostGIS se prefiere a MariaDB por la combinación de transacciones y consultas geográficas del dominio de rutas y sitios; la prueba cubre consultas espaciales y migración.
  \item Keycloak se prefiere a un proveedor de identidad exclusivamente en nube porque permite administrar personal propio y externo sin trasladar la autorización de negocio fuera de M1--M12; el relevo local requiere la prueba de 24 horas.
  \item API Gateway reduce administración frente a un gateway propio, sujeto a prueba de cuotas y costo bajo el pico.
  \item RabbitMQ local más SQS FIFO conserva el trabajo del sitio sin enlace; la deduplicación y el orden por partición se prueban en los consumidores.
  \item Angular se conserva por sus componentes y contratos TypeScript.
\end{itemize}

Los costos, el emplazamiento y la redundancia se comprueban en 4.2 y en el registro ADR.

El núcleo utiliza tecnologías con alternativas de despliegue como Laravel, Angular, PostgreSQL, RabbitMQ y Keycloak. Eso facilita la portabilidad, pero no elimina la dependencia de los servicios administrados de AWS. La reversibilidad documenta contratos, exportación de datos y sustitución de adaptadores, junto con el esfuerzo de migración exigido por RT-03.07, que 4.2 acota por servicio.

El Anexo 4-P reúne versiones de referencia, soporte y criterios de actualización. Las versiones menores se fijan en archivos de bloqueo y SBOM y se promueven mediante pruebas de contrato; no se congela una versión sin soporte por los 56 meses. Laravel, PHP, PostgreSQL, Angular y RabbitMQ se revisan conforme a sus políticas oficiales (Laravel, 2026; PHP, 2026; PostgreSQL, 2026; Angular, 2026a, 2026b; RabbitMQ, 2026).
